\section{Khảo sát liên quan}
Chương này tập trung phân tích và so sánh ưu, nhược điểm của các giải pháp hiện có trên thị trường. Thông qua đó, những đặc tính ưu việt sẽ được chọn lọc nhằm xác định hướng phát triển tối ưu cho đề tài.

\subsection{So sánh các hệ thống hiện tại trên thị trường}

Để có cơ sở thực tiễn cho việc đề xuất và phát triển giải pháp tại thị trường trong nước, chúng tôi đã tiến hành khảo sát, phân tích và đánh giá một số ứng dụng học toán tư duy đang phổ biến. Kết quả so sánh được tổng hợp trong Bảng \ref{tab:math_learning_apps}, tập trung làm rõ sự tương thích giữa các giải pháp quốc tế với đặc thù giáo dục tại Việt Nam.

\begin{table}[h!]
	\centering
	\small
	\renewcommand{\arraystretch}{1.2}
	\resizebox{\textwidth}{!}{%
		\begin{tabular}{|l|p{3.5cm}|p{2.5cm}|p{2cm}|p{2cm}|p{3cm}|p{3.5cm}|}
			\hline
			\textbf{Tên Ứng dụng} & \textbf{Tính năng nổi bật} & \textbf{Đối tượng phù hợp} & \textbf{Ngôn ngữ} & \textbf{Mức phí} & \textbf{Nội dung toán tư duy} & \textbf{Hạn chế tại VN} \\
			\hline
			
			\textbf{Khan Academy Kids} & 
			Video giảng giải chi tiết. \newline 
			Hoạt hình sinh động. & 
			Học sinh có khả năng ngoại ngữ tốt. & 
			Tiếng Anh. & 
			\textbf{Miễn phí}. & 
			Phong phú, chuẩn quốc tế. & 
			Rào cản ngôn ngữ rất lớn; nội dung chưa đúng với khung chương trình Bộ GD\&ĐT Việt Nam. \\
			\hline
			
			\textbf{Duolingo Math} & 
			Game hóa (streak, bảng xếp hạng). & 
			Học sinh muốn học toán kết hợp tiếng Anh. & 
			Tiếng Anh. & 
			Miễn phí cơ bản. & 
			Cơ bản, thiên về luyện tập nhanh. & 
			Chưa hỗ trợ tiếng Việt; thiếu các dạng toán đố, toán tư duy đặc thù của Việt Nam. \\
			\hline
			
			\textbf{Kurio} & 
			Học đa lĩnh vực (Toán, Khoa học). & 
			Học sinh Việt Nam cần kiến thức tổng hợp. & 
			Tiếng Việt. & 
			Có phí. & 
			Nhiều chủ đề nhưng chưa sâu. & 
			Nội dung toán tư duy chưa được chuyên sâu hóa; mức độ tương tác game còn thấp. \\
			\hline
			
			\textbf{VioEdu} & 
			Bám sát SGK Việt Nam. \newline Đấu trường toán học. & 
			Học sinh cần luyện thi và học trên lớp. & 
			Tiếng Việt. & 
			Có phí. & 
			Chủ yếu theo sát chương trình học chính khóa. & 
			Thiên về điểm số và luyện bài tập; thiếu hụt các chuyên đề tư duy logic mở rộng. \\
			\hline
			
			\textbf{Matific} & 
			Học qua mini-game tương tác cao. & 
			Học sinh thích trải nghiệm quốc tế. & 
			Đa ngôn ngữ (hạn chế tiếng Việt). & 
			Phí cao. & 
			Rất sâu sắc về bản chất toán học. & 
			Chi phí đăng ký cao so với thu nhập bình quân tại VN; rào cản về thanh toán quốc tế. \\
			\hline
		\end{tabular}%
	}
	\caption{Bảng so sánh các ứng dụng học toán tư duy cho học sinh tiểu học trong bối cảnh Việt Nam}
	\label{tab:math_learning_apps}
\end{table}

\subsection{Điểm mạnh và điểm yếu của các giải pháp hiện có tại Việt Nam}

Kết quả khảo sát cho thấy bức tranh chung về thị trường ứng dụng giáo dục toán học tại Việt Nam như sau:

\textbf{Điểm mạnh:}
\begin{itemize}
	\item Các ứng dụng nội địa (VioEdu) có ưu thế tuyệt đối về việc bám sát nội dung kiểm tra, đánh giá theo chương trình giáo dục phổ thông của Việt Nam.
	\item Các ứng dụng quốc tế mang lại trải nghiệm người dùng xuất sắc, giúp học sinh tiếp cận toán học dưới góc nhìn trực quan và hiện đại.
\end{itemize}

\textbf{Điểm yếu và hạn chế đặc thù tại Việt Nam:}
\begin{itemize}
	\item \textbf{Sự tách rời giữa "Toán tư duy" và "Toán SGK":} Hiện nay học sinh Việt Nam thường phải chọn một trong hai: hoặc học bám sát chương trình trên lớp nhưng khô khan, hoặc học toán tư duy quốc tế nhưng lại không hỗ trợ được cho các kỳ thi tại trường.
	\item \textbf{Rào cản ngôn ngữ và địa lý:} Các ứng dụng chất lượng hàng đầu thế giới hầu như chưa được Việt hóa hoàn toàn, gây khó khăn cho học sinh ở các vùng tỉnh lẻ hoặc những gia đình chưa chú trọng tiếng Anh sớm.
	\item \textbf{Chi phí tiếp cận:} Các nền tảng học tập chất lượng cao thường có mức giá chưa phù hợp với điều kiện kinh tế của đại đa số gia đình Việt Nam, đặc biệt là ở phân khúc phổ thông.
\end{itemize}

\subsection{Giải pháp đề xuất}

Để giải quyết các bài toán trên, hệ thống học toán tư duy tập trung vào việc \textbf{bản địa hóa trải nghiệm quốc tế} và \textbf{nâng tầm chương trình nội địa}:

\begin{itemize}
	\item \textbf{Nội dung được thiết kế cho học sinh Việt:} Xây dựng hệ thống bài học vừa đảm bảo bám sát chương trình của Bộ Giáo dục và Đào tạo, vừa tích hợp các chuyên đề tư duy logic (Math Logic) theo tiêu chuẩn quốc tế nhưng hoàn toàn bằng tiếng Việt.
	\item \textbf{Game hóa thuần Việt:} Thiết kế giao diện và các thử thách học tập gần gũi với văn hóa, tâm lý của trẻ em Việt Nam, giúp tạo động lực tự học mà không cần sự ép buộc từ phụ huynh.
	\item \textbf{Cá nhân hóa theo năng lực thực tế:} Hệ thống chẩn đoán trình độ dựa trên mặt bằng kiến thức chung của học sinh trong nước, từ đó đưa ra lộ trình bù đắp lỗ hổng kiến thức kịp thời.
	\item \textbf{Hỗ trợ phụ huynh tối đa:} Cung cấp kênh báo cáo tự động qua các nền tảng phổ biến tại Việt Nam, giúp phụ huynh nắm bắt điểm mạnh/yếu của con mình một cách khoa học mà không cần giỏi chuyên môn toán.
	\item \textbf{Khả năng phổ cập:} Hướng tới mục tiêu giúp mọi trẻ em Việt Nam, dù ở thành thị hay nông thôn, đều có cơ hội tiếp cận phương pháp học toán tư duy hiện đại với mức chi phí tối thiểu hoặc miễn phí.
\end{itemize}